\chapter{Conclusion}
\label{ch:conclusion}

The program set out to find a single, neuro-inspired recurrent vision transformer that
learns covert spatial attention from reward, reproduces primate signatures, and generalises
across tasks. After roughly thirty variants and two task horizons, the picture is clearer
than the variant count suggests, because the results collapse onto a few principles.

\textbf{Perception comes first.} The most expensive lesson was the cheapest to fix: sparse
reward cannot teach a front-end to see orientation, and a reconstruction-pretrained VAE
solves it outright. Much of the earlier difficulty on the short task was perception failing
in an architecture's clothes, and separating the two questions --- can it see, and is it
wired right --- is the methodological gain that makes the rest of the comparisons clean.

\textbf{The horizon is a design variable.} The same task at two lengths admits two different
solutions, and recognising that the long task hides perception failures behind temporal
integration explains a paradox that had stalled the program. The short task, with the cue's
information forced through memory and the change forced through perception, is the sharper
instrument.

\textbf{Wiring decides the phenotype.} Coupling value to action through a shared memory is a
precondition for learning, not an optimisation; a grounded image pathway is required; and
which stream drives the policy matters as much as what the streams represent. These routing
facts, not any single clever module, are what separate the models that learn from the models
that collapse.

\textbf{The model already reproduces the paper, and the matched family settled the
mechanism.} With perception fixed, the network reproduces the reliability-scaled cueing
benefit and the causal attention-to-sensitivity link. It also surfaced a genuine mechanistic
question --- a cueing effect carried by memory rather than by attention orienting --- which
the matched family of six wirings has now resolved: \emph{no} variant orients to the cue on
the short task, so the effect is memory-borne in general, not a quirk of one gate. The same
sweep named the network: the two-LSTM multiplicative form is the best detector and a
minimal, paper-faithful choice.

What remains is execution: with the attention mechanism settled and the frozen VAE and
stable harness fixed as constants, run that one architecture across the task battery, and
test whether a longer-delay regime ever restores the published cue-orienting figure. That is
the path from a zoo of variants to the single, interpretable, generalising model the program
was after --- and to an empirical manuscript whose figures, behavioural and neural, are
produced by the same network, with a cueing mechanism reported as it actually is.
